% Level 3, Week 10 -- Project Studio 1
\lesson{10}{Project Studio 1}{Design first: classes, data and milestones -- then build the core.}
\lessoninfo{Concepts: CRC cards, class diagrams, decomposition, milestones, git workflow, code review}{Tools: index cards, Thonny, Terminal + git}{CSTA: 3A-AP-16, 3A-AP-17, 3A-AP-20, 3A-AP-22}

\begin{objectives}
  \item The student can design their project with CRC cards and a class diagram that shows inheritance (\enquote{is a}) and composition (\enquote{has a}). (L3.G1)
  \item The student can split their project into small milestones and choose the data structure and algorithm it needs. (L3.G5, L3.G3)
  \item The student can start a git repository for the project and make small commits with meaningful messages. (L3.G4)
  \item The student can check whether they may use an image, sound or library, and give credit for it. (L3.G9)
\end{objectives}

\begin{materials}
  \item Print: workbook chapter 10 and homework 10. Students bring their pitch card (homework 9).
  \item 10 index cards per student, pencils, sticky notes; a timer.
  \item On each Mac: Terminal and Thonny open. A folder \texttt{Documents/coding-projects} (create it in Finder: \menu{File} $\rightarrow$ \menu{New Folder}).
  \item Prep (10 min): read this plan and the Project teacher guide's \enquote{Milestone 2} page; skim the reference Maze Muncher classes in \texttt{3-advanced/project/reference}.
\end{materials}

\begin{timeplan}
  \timeslot{0--8}{Warm-up}{Pitch round: each student pitches their project in 30 seconds; listeners ask one question each.}
  \timeslot{8--18}{Concept}{Design first: CRC cards, class diagram (is-a / has-a), milestones, small commits.}
  \timeslot{18--28}{Guided demo}{You model CRC cards for Maze Muncher, turn them into a class diagram, then make the first commit in Terminal.}
  \timeslot{28--52}{Hands-on build}{CRC cards (unplugged), class diagram, \texttt{git init}, class skeletons, first commits. Last 6 min: partner code review.}
  \timeslot{52--60}{Wrap-up}{Milestone 2 check, licences talk, exit ticket, homework (project work).}
\end{timeplan}

\section{Warm-up (0--8 min)}
Each student reads their pitch card in 30 seconds; every listener asks one curious question (\enquote{What happens when you win?}). If a pitch is very big, ask: \enquote{What is the smallest version that would still be fun?} Write it on a sticky note: that is Milestone 2's target.

\section{Concept (8--18 min)}
\begin{teacherscript}
\enquote{Builders draw plans before they pour concrete. Programmers do too. Today we \textbf{design first}.

A \textbf{CRC card} is an index card for one class. At the top: the \textbf{Class} name. On the left: its \textbf{Responsibilities} -- what it knows and what it does. On the right: its \textbf{Collaborators} -- the other classes it works with.

Then we draw a \textbf{class diagram}: a box per class, with arrows. A hollow-headed arrow means \emph{is a} -- inheritance: a Chaser \emph{is a} Ghost. A diamond line means \emph{has a} -- composition: a Game \emph{has a} Maze.

Last, we cut the work into \textbf{milestones}, small enough to finish and commit. Commit every time something small works. Your project needs at least eight meaningful commits by the showcase.}
\end{teacherscript}

\begin{analogy}[LEGO instructions]
A big LEGO set comes with numbered bags and a booklet. Bag 1 builds the base, bag 2 the walls. You never open all the bags at once. Milestones are your bags; commits are the photos you take after each bag, so you can always go back.
\end{analogy}

\begin{teachernote}[The CS idea, for you]
Good designs give each class \emph{one} clear job; a long responsibility list means two classes. Inheritance is for \enquote{a kind of}, composition for \enquote{made of}; when unsure, prefer composition. You need not judge whether a design is perfect -- ask the student \emph{why} each class exists.
\end{teachernote}

\section{Guided demo (18--28 min)}
\begin{workedexample}[From cards to a diagram to a first commit]
\textbf{1. CRC cards (on real index cards, 4 min).} Think aloud while writing:
\begin{center}\small
\begin{tabular}{@{}p{2.2cm}p{6.4cm}p{3.6cm}@{}}
\toprule
\textbf{Class} & \textbf{Responsibilities} & \textbf{Collaborators}\\ \midrule
Maze & load the grid from a text file; say if a cell is a wall; count dots & -- \\
Player & know its row and column; move unless a wall is in the way; keep score & Maze \\
Ghost & know its position; move one step (randomly) & Maze \\
Chaser (is a Ghost) & move one step towards the player using breadth-first search with a queue & Maze, Player \\
Game & has a Maze, a Player and Ghosts; runs turns; saves high scores to a file & all of them \\
\bottomrule
\end{tabular}
\end{center}
\textbf{2. Class diagram (2 min).} Draw the boxes as in the workbook diagram: Chaser $\rightarrow$ Ghost with a hollow arrow (\enquote{is a}); Game to Maze, Player, Ghost with diamond lines (\enquote{has a}).

\textbf{3. First commit (4 min).} In Terminal, type and explain each line:
\begin{shellcode}
cd ~/Documents/coding-projects
mkdir maze-muncher
cd maze-muncher
git init
# make muncher_logic.py in Thonny with empty classes, save it here
git add muncher_logic.py
git commit -m "Add empty Maze, Player and Ghost classes"
git log --oneline
\end{shellcode}
Say: \enquote{The message says \emph{what} changed, starting with a verb: Add, Fix, Make, Test.}
\end{workedexample}

\section{Hands-on build (28--52 min)}
\begin{unplugged}[Paper class diagram with CRC cards (8 min)]
Computers closed. Each student writes one index card per class: name at the top, responsibilities left, collaborators right. At least two classes; one must use inheritance or composition. Then \enquote{play} one turn of the game with the cards: hold up each card as it does its job (\enquote{I am Player: I ask Maze if the next cell is a wall}). A card nobody holds up is not needed; a card held up all the time is doing too much. Copy the result into workbook Exercise 10.1.
\end{unplugged}

\begin{enumerate}[itemsep=2pt]
  \item \textbf{Class diagram} (3 min): workbook Exercise 10.2.
  \item \textbf{Repository} (3 min): copy \texttt{3-advanced/project/starter} into \texttt{Documents/coding-projects}, rename it to your project name, then in Terminal: \texttt{cd} into it, \texttt{git init}, first commit.
  \item \textbf{Class skeletons} (5 min): each class with \texttt{\_\_init\_\_} and empty methods (\py|pass|). Commit after each class.
\end{enumerate}

\begin{pairroles}[Architect and builder]
Each student builds their own project, but pairs sit together. While one types, the partner is the \textbf{architect}: holds the CRC cards and asks \enquote{which card does that line belong to?} Swap every 8 minutes. With 1 student: you are the architect. With 3: rotate architect, builder and \textbf{git keeper} (reminds \enquote{commit now!} when something works).
\end{pairroles}

\begin{bughunt}[Partner code review (last 6 minutes)]
Partners swap seats and read each other's code for 3 minutes each, using three questions: (1) Does every class match a CRC card? (2) Is there a method that does two jobs? (3) Is there a name I do not understand? Write one \enquote{star} (something good) and one \enquote{wish} on a sticky note. Common finds: \texttt{self} missing in a method, a class doing everything (\enquote{God class}), \texttt{Ghost.\_\_init\_\_} not called from \texttt{Chaser} (missing \py|super().__init__()|).
\end{bughunt}

\begin{freebuild}[Make it look like yours]
Skeleton committed early? Pick Okabe--Ito colours, name your game, sketch a title screen.
\end{freebuild}

\section{Wrap-up (52--60 min)}
\begin{safety}[Licences: may I use this?]
Say: \enquote{Everything in your project must be yours or allowed. The Python standard library is free to use. Your own drawings and sounds are yours. If you want someone else's picture or sound: look for its \textbf{licence}. \emph{CC0} means free for anyone; \emph{CC BY} means free if you give credit; \enquote{All rights reserved} or no licence means \textbf{no}. Never download anything without your parent, and write every borrowed thing in a Credits section of your README.} Easiest choice: draw everything with turtle.
\end{safety}

\begin{formativecheck}[Milestone 2 check and exit ticket]
\textbf{Look for (M2):} at least two CRC cards; one is-a or has-a link; a git repository with at least 2 commits; class skeletons that run without errors.
\textbf{Question:} \enquote{Which of your classes uses inheritance or composition? Say it as \emph{X is a Y} or \emph{X has a Y}.}
\textbf{Answer:} any correct sentence, e.g. \enquote{Chaser is a Ghost}, \enquote{Gallery has a list of Shapes}.
\end{formativecheck}

\begin{reflection}
Which class was hardest to describe on a card? What is the smallest version of your project that would still be fun? What will your next commit be?
\end{reflection}

Hand out homework 10 (project work: finish the core classes, 3+ commits).

\section{Notes for the teacher}
\begin{commonmistakes}
  \item \textbf{Too big}: a design with 8 classes. Keep 2--4; extra ideas go on a \enquote{later} list.
  \item \textbf{Inheritance for \enquote{has a}}: \texttt{class Game(Maze)}. Ask: \enquote{Is a game a kind of maze?} If not, use composition: \py|self.maze = Maze(...)|.
  \item \textbf{\texttt{git init} in the wrong folder} (e.g. home folder). Check with \texttt{pwd} first; if wrong, delete the hidden \texttt{.git} folder in that folder only with your help.
\end{commonmistakes}

\begin{differentiation}
\supportpath Start from the Maze Muncher design and change the theme (e.g. a bee collecting pollen, avoiding wasps). Two classes are enough: Player and Ghost-like enemy, plus a Game.
\extensionpath Add a module diagram: which file holds which classes (logic file vs drawing file)? Write the first pytest test now.
\end{differentiation}

\begin{stretchchallenge}[Interface first]
Before writing any method body, write every method's name, parameters and a one-line docstring. Can your partner understand the whole design from the docstrings alone?
\end{stretchchallenge}

\begin{teachernote}[If you are short on time]
Shorten the pitch round to 15 seconds each and skip the free-build. Never skip the first commit: students must leave with a repository.
\end{teachernote}
